%    Include other referenced packages here.
\usepackage{pgfplots}
\pgfplotsset{compat=1.15}
\usepackage{mathtools}
\usepackage{amssymb,amsmath,amsthm}
\usepackage[inline]{enumitem}
\usepackage[nosolutionfiles]{answers}
%\usepackage{answers}
%\usepackage{amsrefs}
\usepackage{relsize}
\usepackage{multicol}
\usepackage{booktabs}
\usepackage{hyperref}
\usepackage{calc}
\usepackage{esdiff}
\usepackage{systeme}
\usepackage{bm}
\usepackage{tikz-cd}

% \newtheorem{exercise}{Exercise}
\Newassociation{solution}{extsolution}{solution}

%%%%%%%%%%%%%%%%%%%%%%%%%%%%%%%%%%%%%%%%%%%%%%%%%%%%%%%%%%%%%%%%%%%%%%%%
% definitions
% numerical sets
\newcommand{\real}{{\mathbb R}}
\newcommand{\rn}{{{\mathbb R}^n}}
\newcommand{\realm}{{{\mathbb R}^m}}
\newcommand{\rational}{{\mathbb Q}}
\newcommand{\integer}{{\mathbb Z}}
\newcommand{\naturals}{{\mathbb N}}
\newcommand{\complex}{{\mathbb C}}
\newcommand{\F}{{\mathbb F}}
\newcommand{\R}{{\mathbb R}}
\newcommand{\C}{{\mathbb C}}
\newcommand{\Rn}{{\mathbb R}^n}
\newcommand{\Rm}{{\mathbb R}^m}
\newcommand{\cn}{{\mathbb C}^n}
\newcommand{\rnn}{{{\mathbb R}^{n\times n}}}
\newcommand{\rnm}{{{\mathbb R}^{n\times m}}}
\newcommand{\rmn}{{{\mathbb R}^{m\times n}}}
\newcommand{\cnn}{{{\mathbb C}^{n\times n}}}

% abs and norm, ceil and floor

\DeclarePairedDelimiter\abs{\lvert}{\rvert}%
\DeclarePairedDelimiter\norm{\lVert}{\rVert}%
\DeclarePairedDelimiter\ceil{\lceil}{\rceil}
\DeclarePairedDelimiter\floor{\lfloor}{\rfloor}

% Swap the definition of \abs* and \norm*, so that \abs
% and \norm resizes the size of the brackets, and the 
% starred version does not.
\makeatletter
\let\oldabs\abs
\def\abs{\@ifstar{\oldabs}{\oldabs*}}
%
\let\oldnorm\norm
\def\norm{\@ifstar{\oldnorm}{\oldnorm*}}
\makeatother

% common functions
\newcommand{\rtr}{\real\to\real}
\newcommand{\rttr}{\real^2\to\real}

% common matrices
\newcommand{\MmnF}{M_{mn}(\F)}
\newcommand{\MnF}{M_{n}(\F)}

% diagonal matrix
\DeclareMathOperator{\diag}{diag}

% image of a linear operator
\DeclareMathOperator{\Ima}{Im}

% range of a linear operator
\DeclareMathOperator{\range}{range}

% Rank of a linear operator
\DeclareMathOperator{\rank}{rank}

% projection of a linear operator
\DeclareMathOperator{\proj}{proj}

% Null and Column space of a matrix
\DeclareMathOperator{\rspace}{\mathcal{R}}
\DeclareMathOperator{\nspace}{\mathcal{N}}
\DeclareMathOperator{\cspace}{\mathcal{C}}

% Trace of a linear operator
\DeclareMathOperator{\tr}{tr}

% Rank of a linear operator
\DeclareMathOperator{\vspan}{span}

% Sign of a function
\DeclareMathOperator{\sign}{sgn}

% differential d
\newcommand{\dif}{\mathop{}\!\mathrm{d}}
\newcommand{\od}[2]{\frac{\dif #1}{\dif #2}}
\newcommand{\dod}[2]{\displaystyle\frac{\dif #1}{\dif #2}}

% complement
\newcommand{\scomp}[1]{{#1^\complement}}

% interior & closure
\newcommand{\interior}[1]{{\kern0pt#1}^{\mathrm{o}}}
\newcommand{\closure}[1]{{\overline{#1}}}

% evaluated integral 
\newcommand\eval[3]{\left[#1\right]_{#2}^{#3}}

% floor and ceiling
%\newcommand\floor[1]{\lfloor#1\rfloor}
%\newcommand\ceil[1]{\lceil#1\rceil}

% base n numbers
\newcommand{\basen}[2]{{#1}_{#2}}

% matrices p q
\newcommand{\rmat}[2]{{{\mathbb R}^{{#1}\times {#2}}}}

% vectors
\newcommand{\vect}[1]{\mathbf{#1}}
\newcommand{\va}{\mathbf{a}}
\newcommand{\vb}{\mathbf{b}}
\newcommand{\vc}{\mathbf{c}}
\newcommand{\vd}{\mathbf{d}}
\newcommand{\ve}{\mathbf{e}}
\newcommand{\vf}{\mathbf{f}}
\newcommand{\vp}{\mathbf{p}}
\newcommand{\vq}{\mathbf{q}}
\newcommand{\vr}{\mathbf{r}}
\newcommand{\vs}{\mathbf{s}}
\newcommand{\vx}{\mathbf{x}}
\newcommand{\vy}{\mathbf{y}}
\newcommand{\vz}{\mathbf{z}}
\newcommand{\vu}{\mathbf{u}}
\newcommand{\vw}{\mathbf{w}}
\newcommand{\vv}{\mathbf{v}}
\newcommand{\vT}{\mathbf{T}}
\newcommand{\vbeta}{\mathbf{\beta}}
\newcommand{\vzero}{\mathbf{0}}
\newcommand{\mzero}{\mathbf{O}}

\newcommand{\vtn}[4]{
\vect{#1}=\left[
\begin{array}{c}
#2 \\
#3 \\
#4
\end{array}
\right]}

\newcommand{\vtu}[3]{
\left[
\begin{array}{c}
#1 \\
#2 \\
#3
\end{array}
\right]}

\newcommand{\vtd}[4]{
\left[
\begin{array}{c}
#1 \\
#2 \\
#3 \\
#4
\end{array}
\right]}

\newcommand{\vdn}[3]{
\vect{#1}=\left[
\begin{array}{c}
#2 \\
#3 
\end{array}
\right]}

\newcommand{\vdu}[2]{
\left[
\begin{array}{c}
#1 \\
#2 
\end{array}
\right]}

\newcommand{\stlb}[3]{
\left\{\begin{matrix}
#1 \\ #2 \\ #3
\end{matrix}\right.}

\newcommand{\stlt}[3]{
\left\{\begin{matrix}
#1 \\ #2 \\ #3
\end{matrix}\right.}

% horbars in matrices
\newcommand{\brows}[1]{%
  \begin{bmatrix}
  \begin{array}{@{\protect\rotvert\;}c@{\;\protect\rotvert}}
  #1
  \end{array}
  \end{bmatrix}
}
\newcommand{\rotvert}{\rotatebox[origin=c]{90}{$\vert$}}
\newcommand{\rowsvdots}{\multicolumn{1}{@{}c@{}}{\vdots}}

% tikz
\usepackage{tikz,ifthen,calc}
\newcounter{mysum}
\newcounter{myline}
\usetikzlibrary{shapes,arrows,shadows,trees,patterns}
\usetikzlibrary{fadings}
\tikzfading[name=fade out, inner color=transparent!0, outer color=transparent!100]

\tikzstyle{2dgraph}=[
			interior1/.style={red!20!white},
			boundary1/.style={very thick,red},
			interior2/.style={blue!30!white},
			boundary2/.style={very thick,blue},
			boundary2b/.style={dashed,blue},
			interior12/.style={orange!10!cyan!10!white},
			obj/.style={very thick,blue},
			objval/.style={sloped,midway,above,font=\footnotesize},
			bcircle/.style={rectangle, rounded corners,draw,green!80!white,very thick,minimum height=0.6cm},
			axes/.style={->,thick,black}]
\usetikzlibrary{decorations.pathreplacing,decorations.pathmorphing,decorations.markings,matrix}

\usepgfplotslibrary{fillbetween}

\tikzstyle arrowstyle=[scale=1]
\tikzstyle directed=[postaction={decorate,decoration={markings,
    mark=at position .65 with {\arrow[arrowstyle]{stealth}}}}]

% Venn diagrams

\def\firstcircle{(0,0) circle (1.5cm)}
\def\secondcircle{(0:2cm) circle (1.5cm)}
\def\universe{(-2,2) rectangle (4,-2)}

\colorlet{circle edge}{blue!50}
\colorlet{circle area}{blue!20}
\colorlet{empty area}{white}

\tikzset{
	filled/.style={fill=circle area, draw=circle edge, thick},
	nofilled/.style={fill=empty area, draw=circle edge, thick},
    outline/.style={draw=circle edge, thick}}

\tikzset{
	invisible/.style={opacity=0},
	visible on/.style={alt={#1{}{invisible}}},
	alt/.code args={<#1>#2#3}{%
	\alt<#1>{\pgfkeysalso{#2}}{\pgfkeysalso{#3}} % \pgfkeysalso doesn't change the path
	},
}

%%%%%%%%%%%%%%%%%%%%%%%%%%%%%%%%%%%%%%%%%%%%%%%%%%%%%%%%%%%%%%%%%%%%%%%%
% theorems
%\newtheorem{theorem}{Theorem}[chapter]
%\newtheorem{lemma}[theorem]{Lemma}
%
%\theoremstyle{definition}
%\newtheorem{definitionx}[theorem]{Definition}
%\newenvironment{definition}
%  {\pushQED{\qed}\renewcommand{\qedsymbol}{$\Diamond$}\definitionx}
%  {\popQED\enddefinitionx}
%
%% \newtheorem{example}[theorem]{Example}
%\newtheorem{examplex}[theorem]{Example}
%\newenvironment{example}
%  {\pushQED{\qed}\renewcommand{\qedsymbol}{$\triangle$}\examplex}
%  {\popQED\endexamplex}
%  
\newtheorem{exercise}{Exercise}[]
%
%\theoremstyle{remark}
%\newtheorem{remark}[theorem]{Remark}
%
%\numberwithin{section}{chapter}
%\numberwithin{equation}{chapter}


%
% for date in the title
%
\makeatletter
\def\@adminfootnotes{%
  \let\@makefnmark\relax  \let\@thefnmark\relax
  \ifx\@empty\@date\else \@footnotetext{\@setdate}\fi%%   <------ added
  \ifx\@empty\@subjclass\else \@footnotetext{\@setsubjclass}\fi
  \ifx\@empty\@keywords\else \@footnotetext{\@setkeywords}\fi
  \ifx\@empty\thankses\else \@footnotetext{%
    \def\par{\let\par\@par}\@setthanks}%
  \fi
}
\makeatother